% Tarea 1 - Informe de entrevista
\documentclass[12pt,a4paper]{article}
\usepackage[utf8]{inputenc}
\usepackage[T1]{fontenc}
\usepackage[spanish]{babel}
\usepackage{geometry}
\usepackage{parskip}
\geometry{left=2.5cm,right=2.5cm,top=2.5cm,bottom=2.5cm}

\title{Tarea 1 -- Informe de entrevista sobre comunicaciones}
\author{John Ureña \\
ID: 1085857 \\
Instituto Tecnológico de Santo Domingo}
\date{}

\begin{document}
\maketitle

\textbf{Entrevistado:} Juan Fulvio Ureña Meléndez (padre)

\section*{Introducción}
Para esta tarea entrevisté a mi padre, Juan Fulvio Ureña Meléndez, con el objetivo de conocer cómo vivió de primera mano la evolución de las tecnologías de comunicación desde su niñez hasta la actualidad. A partir de sus respuestas y de una breve investigación complementaria sobre las mejoras que traerá la tecnología 5G, elaboré el siguiente informe.

\section*{Desarrollo de la entrevista}
Al preguntarle cuál era el medio de comunicación más utilizado en su hogar durante su niñez, mi padre respondió sin dudar que, en aquella época, la radio era el medio predominante en la casa. Era el aparato al que todos recurrían para enterarse de las noticias y para el entretenimiento familiar.

Le pregunté también qué diferencias notaba entre la tecnología de los años 90 y la que usamos hoy en día, y su respuesta trajo consigo un recuerdo bastante gracioso: su primer teléfono celular era tan grande que lo describió como "un ladrillo", con todo y antena retráctil. Según me contó, en términos generales los equipos de esa época eran mucho más voluminosos y menos prácticos que los actuales, además de tener una capacidad de procesamiento y de conectividad muy limitada en comparación con lo que tenemos hoy.

Sobre las dificultades que enfrentó al usar Internet en sus primeros años, mencionó que, más allá del propio choque cultural que significó adaptarse a esa nueva tecnología, la conexión era lenta e inestable. Recuerda con frecuencia los cortes de conexión y la baja velocidad, algo que hacía que navegar fuera una experiencia bastante distinta a la fluidez con la que contamos actualmente.

Finalmente, le pedí que compartiera alguna anécdota nostálgica relacionada con las telecomunicaciones, y me contó que de niño tenía que ajustar manualmente la antena de la televisión para poder sintonizar bien la señal y así poder ver la lucha libre sin que la imagen se perdiera. Este tipo de detalles cotidianos reflejan bastante bien cómo era vivir la tecnología de comunicaciones en esa época.

\section*{Investigación: mejoras que traerá la implementación del 5G}
Complementando la entrevista, investigué un poco sobre qué cambios traerá consigo la implementación de la tecnología 5G. En términos generales, se espera que esta tecnología impacte de forma significativa a distintos sectores productivos. Uno de los beneficios más mencionados es el aumento en la velocidad de conexión junto con una reducción notable en la latencia, lo cual abre la puerta a comunicaciones en tiempo real para aplicaciones críticas como la telemedicina o el control remoto de procesos industriales.

Otro punto importante es la capacidad del 5G para soportar la conexión masiva de dispositivos al mismo tiempo, lo que resulta fundamental para el crecimiento del Internet de las cosas (IoT) en áreas como las ciudades inteligentes, la agricultura y la logística. A esto se suma que su baja latencia y alta fiabilidad la hacen una tecnología ideal para aplicaciones como los vehículos autónomos o el uso de realidad aumentada y virtual en la educación y el entretenimiento.

\end{document}
